\begin{answers}{ch:lora}{LoRA：少ないパラメータで適応させる}
\item ランクは 1 で、$\begin{pmatrix}1&2\\3&6\\-1&-2\end{pmatrix}=\begin{pmatrix}1\\3\\-1\end{pmatrix}\begin{pmatrix}1&2\end{pmatrix}$ と分解できます。
2 列目は 1 列目のちょうど 2 倍なので、列ベクトルの向きは 1 つだけです。
行で見ても、2 行目は 1 行目の 3 倍、3 行目は $-1$ 倍です。
1 行目 $(1,2)$ を共通の行ベクトルにし、各行の倍率 $(1,3,-1)$ を列ベクトルにすれば、積が元の行列になります。
\item 条件は $r\le d/200$ で、$d=4096$ なら $r\le20$ です。
元の重みは $d^2$ 個、LoRA は\eqref{eq:lora-count}より $r(d+d)=2rd$ 個です。
$2rd\le0.01\,d^2$ の両辺を $2d$ で割ると $r\le d/200$ です。
$4096/200=20.48$ で、$r$ は整数なので 20 以下です（$r=16$ なら約 0.78\%）。
\item 両方が $\vzero$ だと、\eqref{eq:lora-grad-b}は $\mA\vx=\vzero$ のため $\vzero$、\eqref{eq:lora-grad-a}も $\mB=\vzero$ のため $\vzero$ になります。
勾配が 0 なので、何回更新しても $\mA$ と $\mB$ は $\vzero$ のまま動きません。
片方（本文では $\mA$）を乱数にしておけば、$\mB$ の勾配が 0 でなくなり、学習が始まります。
逆に両方を乱数にすると、学習を始める前から $\frac{\alpha}{r}\mB\mA\ne\vzero$ です。
出発点が事前学習の重み $\mW_0$ からずれてしまい、せっかく学んだ能力を、でたらめな差分で壊した状態から学習することになります。
\item 本文の実験の出力（\file{ch17/out/lora_r8_qkv.txt}、\file{lora_r8_mlp.txt}）から答えます。
$r=8$ で \code{qkv} だけなら 24,576 個で正答率 8.6\%、MLP だけなら 61,440 個で 68.1\% で、\tabref{tab:lora-targets}の値そのものです（全線形層は 98,304 個で 81.7\%）。
同じパラメータ数で比べるには、\eqref{eq:lora-count}の $r(d_{\text{in}}+d_{\text{out}})$ を全部の層で足した数がそろうように $r$ を選びます。
パラメータ数は $r$ に比例するので、$r=8$ の数を 8 で割れば $r=1$ あたりの数が出ます（全線形層 12,288、\code{qkv} 3,072、MLP 7,680）。
\code{qkv} の $r=8$ とそろえるなら全線形層で $r=2$（24,576 個。本文では 68.6\%）、MLP の $r=8$ とそろえるなら全線形層で $r=5$（61,440 個）です。
$r=5$ は実行していませんが、全線形層の $r=4$（79.4\%）と $r=8$（81.7\%）の間に入り、MLP だけの 68.1\% を上回ると予想されます。
\item 全パラメータ更新では約 280 GB、LoRA では約 15.6 GB で、LoRA は約 18 分の 1 です。
\tabref{tab:lora-cost}より、全パラメータ更新は用途ごとに 14 GB のコピーが要るので、$14\times20=280$ GB です。
LoRA は共通の基盤モデル 14 GB を 1 つだけ持ち、用途ごとに 80 MB の差分を足すので、$14+0.08\times20=15.6$ GB です。
用途を増やしたときの増え方も、1 つあたり 14 GB と 80 MB で、約 175 倍違います。
推論のときも、基盤モデルを 1 つメモリに載せたまま、差分を差し替えるだけで用途を切り替えられます。
\end{answers}
